\section{Weakly Supervised Learning}
\label{sec:weakly}

Every classifier so far was trained on a dataset in which each example carries
a correct and precise label. Such datasets are expensive. A pixel-accurate
segmentation mask takes minutes per image, medical labels need experts, and
labels collected cheaply from web tags or crowd workers are often wrong. In
practice one usually has something weaker: many unlabeled images and only a
few labeled ones, labels for a whole image but not for the object inside it,
or labels of which some fraction is simply wrong. \emph{Weakly supervised
learning} (WSL) collects the methods that still train useful models from such
data. This section first sorts weak supervision into three types, then treats
each in turn: unlabeled data in active, semi-supervised, and transductive
learning; coarse labels in multi-instance learning; and wrong labels, where
early-learning regularization and negative learning prevent a network from
memorizing the errors. It ends with a critical look at how much weak
supervision actually gains compared with a few clean labels.

\subsection{Three Types of Weak Supervision}

Weak supervision can fail to be ``strong'' in three different ways.
Figure~\ref{fig:weak-types} shows them side by side for a binary problem with
labels \textsf{Y} and \textsf{N}; each bar stands for one feature vector.

\begin{itemize}
  \item \emph{Incomplete supervision}: only a small subset of the training
    data is labeled; the rest has no label at all. Example: $1000$ labeled
    and $100\,000$ unlabeled images of skin lesions.
  \item \emph{Inexact supervision}: labels exist but are coarser than the
    task requires. Example: the image is tagged ``tiger,'' but the task is to
    find the tiger pixels or patches.
  \item \emph{Inaccurate supervision}: every example has a label of the
    desired granularity, but some labels are wrong. Example: web images
    labeled by the surrounding text.
\end{itemize}

\begin{figure}[ht]
  \centering
  \begin{tikzpicture}[
    feat/.style={draw, fill=gray!15, minimum width=1.0cm, minimum height=0.28cm,
      inner sep=0pt},
    lab/.style={draw, circle, minimum size=0.42cm, inner sep=0pt,
      font=\scriptsize\sffamily},
    bag/.style={draw, rounded corners=3pt, thick, gray!70!black},
    ttl/.style={font=\small, align=center}
  ]
    % (a) incomplete
    \foreach \y/\lv/\c in {0/Y/red!15, 0.55//white, 1.1/N/red!15,
        1.65//white, 2.2//white} {
      \node[feat] (f) at (0,\y) {};
      \node[lab, right=0.25cm of f, fill=\c] {\lv};
    }
    \node[ttl] at (0.45,-0.75) {(a) incomplete:\\few labels};
    % (b) inexact
    \begin{scope}[xshift=4.2cm]
      \foreach \y/\lv in {0/N, 1.35/Y} {
        \node[feat] at (-0.1,\y) {};
        \node[feat] at (-0.1,\y+0.4) {};
        \node[feat] at (-0.1,\y+0.8) {};
        \draw[bag] (-0.75,\y-0.28) rectangle (0.55,\y+1.08);
        \node[lab, fill=red!15] at (1.05,\y+0.4) {\lv};
      }
      \node[ttl] at (0.2,-0.75) {(b) inexact:\\one label per bag};
    \end{scope}
    % (c) inaccurate
    \begin{scope}[xshift=8.2cm]
      \foreach \y/\lv/\c in {0/Y/red!15, 0.55/N/red!15, 1.1/N/orange!45,
          1.65/Y/red!15, 2.2/Y/orange!45} {
        \node[feat] (f) at (0,\y) {};
        \node[lab, right=0.25cm of f, fill=\c] {\lv};
      }
      \node[ttl] at (0.45,-0.75) {(c) inaccurate:\\some labels wrong};
    \end{scope}
  \end{tikzpicture}
  \caption{The three types of weak supervision. Bars are feature vectors,
  circles their labels. (a) Most examples have no label. (b) Labels belong
  to bags of instances, not to the instances. (c) Every example has a label,
  but some are wrong (orange); the learner does not know which.}
  \label{fig:weak-types}
\end{figure}

\noindent
The types call for different tools, summarized in
Table~\ref{tab:weak-types}. Incomplete supervision asks how unlabeled data
can help at all. Inexact supervision asks how a label for a group transfers
to its members. Inaccurate supervision asks how to stop a network from
learning the wrong labels. Real datasets often combine the types, e.g.\ a few
noisy image-level tags and many untagged images.

\begin{table}[ht]
  \centering
  \begin{tabular}{@{}p{2.2cm}p{4.4cm}p{3.9cm}p{4.2cm}@{}}
    \toprule
    Type & What is given & Central question & Methods here \\
    \midrule
    Incomplete & few labeled, many unlabeled examples & what do unlabeled
      data reveal? & active learning, semi-supervised and transductive
      learning \\
    Inexact & labels for bags (e.g.\ images), task on instances (patches) &
      which instance caused the label? & multi-instance learning \\
    Inaccurate & a label for every example, some wrong & how to avoid
      learning the errors? & early-learning regularization, negative
      learning, SelNLPL \\
    \bottomrule
  \end{tabular}
  \caption{Taxonomy of weak supervision. Each type has its own question and
  its own family of methods.}
  \label{tab:weak-types}
\end{table}

\subsection{Incomplete Supervision: Learning with Unlabeled Data}

With incomplete supervision, the training set consists of $n$ labeled
examples $\{(x_i,y_i): i=1,\dots,n\}$ and $m$ unlabeled examples
$x'_1,\dots,x'_m$, usually with $m\gg n$. There are two ways to use the
unlabeled part. \emph{Active learning} assumes an \emph{oracle}, typically a
human expert, who can be asked for the true label of selected unlabeled
examples. The learner picks the examples whose labels would help most, for
instance those closest to the current decision boundary, and asks only for
those. The goal is a good model with as few queries as possible.
\emph{Semi-supervised learning} uses the unlabeled data automatically,
without any human in the loop. It comes in two variants that differ in what
the unlabeled data are and what the output is.

\begin{itemize}
  \item \emph{Pure semi-supervised learning} works under an
    \emph{open-world} assumption. The test data are unknown at training time;
    the unlabeled data are not necessarily the test data. The output is a
    classifier $f$ that can label any future input.
  \item \emph{Transductive learning} works under a \emph{closed-world}
    assumption. The test data are given in advance, and they are exactly the
    unlabeled examples. The goal is to label them as well as possible; no
    general function $f$ is needed.
\end{itemize}

\noindent
Figure~\ref{fig:ssl-settings} compares the three workflows. In all three, a
model is trained from labeled and unlabeled data. They differ in whether a
human adds labels during training and in what the model predicts afterwards.

\begin{figure}[ht]
  \centering
  \begin{tikzpicture}[
    db/.style={draw, cylinder, shape border rotate=90, aspect=0.25,
      minimum width=1.25cm, minimum height=0.9cm, align=center,
      font=\scriptsize},
    box/.style={draw, rounded corners=2pt, minimum height=0.8cm,
      minimum width=1.3cm, align=center, font=\scriptsize},
    arr/.style={-{Latex[length=2mm]}, thick},
    ttl/.style={font=\small\bfseries, anchor=west}
  ]
    % row 1: active learning
    \node[ttl] at (-0.8,0.95) {Active learning};
    \node[db] (l1) at (0,0) {labeled};
    \node[db, draw=red!70!black] (u1) at (1.9,0) {unlabeled};
    \node[box, fill=yellow!20] (o1) at (4.3,0) {oracle\\(expert)};
    \node[box] (m1) at (6.9,0) {model};
    \node[box] (t1) at (9.6,0) {unseen\\test data};
    \draw[arr] (u1.east) -- node[above, font=\scriptsize] {query} (o1.west);
    \draw[arr] (o1.south) to[bend left=25]
      node[below, font=\scriptsize, pos=0.5] {true labels} (l1.south);
    \draw[arr] (o1.east) -- node[above, font=\scriptsize] {train} (m1.west);
    \draw[arr] (m1.east) -- node[above, font=\scriptsize] {predict} (t1.west);
    % row 2: pure semi-supervised
    \node[ttl] at (-0.8,-1.75) {Pure semi-supervised learning};
    \node[db] (l2) at (0,-2.7) {labeled};
    \node[db, draw=red!70!black] (u2) at (1.9,-2.7) {unlabeled};
    \node[box] (m2) at (6.9,-2.7) {model};
    \node[box] (t2) at (9.6,-2.7) {unseen\\test data};
    \draw[arr] (u2.east) -- node[above, font=\scriptsize] {train} (m2.west);
    \draw[arr] (m2.east) -- node[above, font=\scriptsize] {predict} (t2.west);
    % row 3: transductive
    \node[ttl] at (-0.8,-3.95) {Transductive learning};
    \node[db] (l3) at (0,-4.9) {labeled};
    \node[db, draw=red!70!black] (u3) at (1.9,-4.9) {unlabeled};
    \node[box] (m3) at (6.9,-4.9) {model};
    \node[db, draw=red!70!black] (t3) at (9.6,-4.9) {same\\unlabeled};
    \draw[arr] (u3.east) -- node[above, font=\scriptsize] {train} (m3.west);
    \draw[arr] (m3.east) -- node[above, font=\scriptsize] {predict} (t3.west);
  \end{tikzpicture}
  \caption{Three ways to use unlabeled data. Active learning asks an oracle
  for labels of selected examples. Pure semi-supervised learning uses the
  unlabeled data without a human and predicts on unseen test data.
  Transductive learning predicts exactly on the unlabeled data it was trained
  with.}
  \label{fig:ssl-settings}
\end{figure}

\paragraph{Induction versus transduction.}
The difference between the settings is clearest when written as mappings.
Ordinary supervised learning is \emph{induction}: from labeled examples, infer
a general function,
\begin{equation}
  \{(x_i,y_i): i=1,\dots,n\} \;\mapsto\; f .
\end{equation}
Pure semi-supervised learning is induction with unlabeled data,
\begin{equation}
  \{(x_i,y_i): i=1,\dots,n\}\cup\{x'_1,\dots,x'_m\} \;\mapsto\; f ,
\end{equation}
and \emph{transduction} skips the function and outputs labels directly,
\begin{equation}
  \{(x_i,y_i): i=1,\dots,n\}\cup\{x'_1,\dots,x'_m\} \;\mapsto\;
  (y'_1,\dots,y'_m).
  \label{eq:transduction}
\end{equation}
The output of transduction is a vector of labels, not a function. The idea,
due to Vapnik, is that one should not solve a harder problem (find $f$ for
every possible input) as an intermediate step when only a simpler one ($m$
particular labels) is asked. Information is transferred directly from the
labeled to the unlabeled examples. The price is that a new test image
requires a new transductive run. Which setting to choose depends on whether
unlabeled data are available, whether an interpretable function $f$ is
needed, and whether one can afford to retrain for every new batch of test
data.

\paragraph{Why unlabeled data help: the cluster assumption.}
An unlabeled example carries no label, so it is not obvious why it helps. It
does carry information about \emph{where the data lie}.
Figure~\ref{fig:cluster-assumption} gives the standard example. With one
labeled point of each class, a test point halfway between them cannot be
decided. After observing unlabeled points, it becomes clear that the test
point lies in the same dense cluster as the positive example. The
\emph{cluster assumption} states that points in the same cluster tend to
share a label, or equivalently that the decision boundary should pass
through low-density regions. Semi-supervised methods rely on this or a
related assumption. If it fails, for example because the classes overlap in
one dense blob, unlabeled data do not help and can even hurt.

\begin{figure}[ht]
  \centering
  \begin{tikzpicture}[
    pt/.style={circle, fill=gray!55, inner sep=0pt, minimum size=5pt},
    lbl/.style={font=\small, align=center}
  ]
    % left: labeled only
    \node[font=\Large] at (0,0) {$+$};
    \node[font=\Large] at (3.0,0) {$-$};
    \draw[thick] (1.2,0) circle (4pt);
    \node[lbl] at (1.2,0.75) {test point:\\$+$ or $-$?};
    % arrow
    \draw[-{Latex[length=2.5mm]}, thick] (3.7,0) -- (5.0,0)
      node[midway, above, font=\scriptsize, align=center]
      {observe\\unlabeled data};
    % right: with unlabeled cluster
    \begin{scope}[xshift=5.8cm]
      \foreach \p in {(0.3,0.55), (0.8,0.8), (1.3,0.6), (0.6,-0.5),
          (1.1,-0.7), (1.6,-0.35), (0.25,-0.1), (1.75,0.25), (1.5,-0.9),
          (2.8,0.7), (3.3,0.45), (3.7,0.8), (3.1,-0.3), (3.6,-0.6),
          (4.0,-0.1), (3.4,-0.95)}
        \node[pt] at \p {};
      \node[font=\Large] at (0,0) {$+$};
      \node[font=\Large] at (3.0,0) {$-$};
      \draw[thick] (1.2,0) circle (4pt);
      \draw[dashed, blue!70!black, thick] (2.25,-1.2) -- (2.25,1.2);
      \node[lbl, text=blue!70!black] at (2.25,1.5) {boundary in gap};
      \node[lbl] at (1.2,-1.45) {test point: $+$};
    \end{scope}
  \end{tikzpicture}
  \caption{The cluster assumption. With two labeled points alone, the test
  point (circle) is ambiguous. Unlabeled points (gray) reveal two clusters;
  the boundary should run through the empty gap, so the test point belongs
  to the positive cluster.}
  \label{fig:cluster-assumption}
\end{figure}

\paragraph{A mixture model with unlabeled data.}
The cluster assumption can be made precise with a mixture model. Assume the
data come from a mixture of components $g=1,\dots,K$ with weights
$\alpha_j$ and densities $f(x\mid\theta_j)$, for example Gaussians. A binary
classifier with classes $c\in\{\textsf{Y},\textsf{N}\}$ decides via the
component that probably generated $x$:
\begin{align}
  h(x) &= \arg\max_{c\in\{\textsf{Y},\textsf{N}\}}
    \sum_{j=1}^{K} P(y=c \mid g=j, x)\; P(g=j\mid x),
  \label{eq:ssl-mixture}\\
  P(g=j\mid x) &= \frac{\alpha_j f(x\mid\theta_j)}
    {\sum_{k=1}^{K}\alpha_k f(x\mid\theta_k)} .
  \nonumber
\end{align}
The second factor, the posterior of component $j$, depends only on $x$. It
can therefore be estimated from \emph{all} examples, labeled or not, e.g.\
with the EM algorithm of Appendix~\ref{app:gmm}. Only the first factor, which
component corresponds to which class, requires labels. If the clusters are
well separated, a handful of labeled examples per component suffices to fix
this assignment, while the shape and position of the components come from
the unlabeled majority.

\paragraph{Worked example.}
Two Gaussian components have been fitted with EM on $10\,000$ images, of which
only $20$ are labeled. Of the labeled images that fall into component $1$
(highest $P(g=1\mid x)$), $9$ are \textsf{Y} and $1$ is \textsf{N}; of those
in component $2$, $1$ is \textsf{Y} and $9$ are \textsf{N}. Then
$P(y=\textsf{Y}\mid g=1)\approx0.9$ and $P(y=\textsf{Y}\mid g=2)\approx0.1$.
A new image with $P(g=1\mid x)=0.7$ gets
\begin{equation*}
  P(\textsf{Y}\mid x)\approx 0.9\cdot0.7 + 0.1\cdot0.3 = 0.66,
  \qquad
  P(\textsf{N}\mid x)\approx 0.1\cdot0.7 + 0.9\cdot0.3 = 0.34,
\end{equation*}
so $h(x)=\textsf{Y}$. Twenty labels could not have estimated two
high-dimensional Gaussians, but they suffice to name the components.

\paragraph{Measuring label efficiency.}
A frequent practical question is how many labels a representation needs. A
\emph{learning curve} answers it: train the same classifier, e.g.\ a linear
SVM, on stratified fractions of the labeled set ($1\,\%$, $5\,\%$,
$10\,\%$, $25\,\%$, $50\,\%$, $100\,\%$) and plot test accuracy against the
fraction. Stratification keeps the class proportions equal in every
fraction, so small fractions are not missing whole classes, and the SVM
parameter $C$ should be re-chosen by cross-validation at each fraction
because the best regularization depends on the amount of data. Comparing a
handcrafted feature such as HOG with frozen self-supervised features such as
DINO (Appendix~\ref{app:foundation}) on the same curve shows directly how
much the unlabeled pretraining saves: a good representation reaches a given
accuracy with a much smaller fraction. This is a form of exploiting unlabeled
data too, only the unlabeled data were used before, during pretraining.

\subsection{Inexact Supervision: Multi-Instance Learning}

In inexact supervision some label information is present, but not at the
level the task needs. The best-known case is \emph{multi-instance learning}
(MIL). The training data are organized into \emph{bags}, i.e.\ sets of
instances, and only the bag carries a label. In the binary case the rule is:
a bag is positive if at least one of its instances is positive, and negative
if all of its instances are negative.

Images fit this pattern naturally, as Figure~\ref{fig:mil} shows. An image is
a bag; its patches are the instances. The label ``tiger'' says that some
patches show the tiger, but not which ones; most patches show grass, trees,
or background. An image without a tiger guarantees that none of its patches
show a tiger. Formally, with bag $B=\{x_1,\dots,x_K\}$, hidden instance labels
$y_k\in\{0,1\}$, and bag label $Y$,
\begin{equation}
  Y = \max_{k=1,\dots,K} y_k .
  \label{eq:mil}
\end{equation}
The goal is to predict labels of unseen bags, and often, as a by-product, to
find the positive instances inside them.

\begin{figure}[ht]
  \centering
  \begin{tikzpicture}[
    cell/.style={draw, minimum size=0.55cm, inner sep=0pt},
    lbl/.style={font=\small, align=center}
  ]
    % positive bag
    \foreach \i in {0,...,4} \foreach \j in {0,...,3}
      \node[cell, fill=green!8] at (\i*0.55,\j*0.55) {};
    \foreach \i/\j in {1/1, 2/1, 2/2, 1/2, 3/2}
      \node[cell, fill=orange!60] at (\i*0.55,\j*0.55) {};
    \draw[very thick] (-0.275,-0.275) rectangle (2.475,1.925);
    \node[lbl] at (1.1,-0.75) {bag label $Y=1$ (``tiger'')};
    \node[lbl, font=\scriptsize] at (1.1,2.35)
      {some instances positive (orange)};
    % negative bag
    \begin{scope}[xshift=5.2cm]
      \foreach \i in {0,...,4} \foreach \j in {0,...,3}
        \node[cell, fill=green!8] at (\i*0.55,\j*0.55) {};
      \draw[very thick] (-0.275,-0.275) rectangle (2.475,1.925);
      \node[lbl] at (1.1,-0.75) {bag label $Y=0$};
      \node[lbl, font=\scriptsize] at (1.1,2.35)
        {all instances negative};
    \end{scope}
    % pooling sketch
    \begin{scope}[xshift=9.6cm]
      \node[draw, rounded corners=2pt, align=center, font=\scriptsize,
        minimum width=1.6cm] (s) at (0,1.3) {patch scores\\$s(x_1),\dots,s(x_K)$};
      \node[draw, rounded corners=2pt, align=center, font=\scriptsize,
        minimum width=1.6cm] (p) at (0,0) {$\max_k s(x_k)$};
      \draw[-{Latex[length=2mm]}, thick] (s) -- (p);
      \node[lbl] at (0,-0.75) {bag score};
    \end{scope}
  \end{tikzpicture}
  \caption{Multi-instance learning on images. An image is a bag of patches.
  A positive image contains at least one positive patch, whose position is
  unknown; a negative image contains none. A network can score every patch
  and pool the scores with a maximum to obtain the bag score.}
  \label{fig:mil}
\end{figure}

\noindent
The asymmetry of Equation~\ref{eq:mil} matters for learning. Every instance
of a negative bag is a certain negative example, so negative bags provide
clean instance-level labels. A positive bag says only that the maximum is
positive, so it constrains just its most positive instance. A common neural
implementation scores each patch with a shared network $s(x_k)$, pools the
scores with a maximum (or a smooth version such as log-sum-exp or attention
pooling), and trains the pooled bag score with the binary cross-entropy on
$Y$. The gradient then flows mainly into the patch that currently has the
highest score. After training, the patch scores themselves act as a coarse
localization, although the image label never said where the object was.

\paragraph{Worked example.}
A bag of four patches receives scores $0.1,\,0.8,\,0.3,\,0.2$ after the
sigmoid. With max pooling the bag probability is $0.8$. If the bag label is
$Y=1$, the loss is $-\log0.8\approx0.22$ and only the second patch receives a
gradient, pushing its score further up. If the bag label is $Y=0$, the loss
is $-\log(1-0.8)\approx1.61$, again concentrated on the second patch, now
pushing it down. Repeated over many bags, patches that appear only in positive
images are pushed up and patches that also appear in negative images are
pushed down.

\subsection{Inaccurate Supervision and Memorization}

In inaccurate supervision every training example has a label, but some
labels are wrong. The learner does not know which ones. Two noise models are
common in experiments. With \emph{symmetric noise}, a fraction of the labels
is replaced by a class drawn uniformly at random. In the
\emph{symmetric-inclusive} variant the random draw includes the true class,
so with $30\,\%$ noise and $C=10$ classes, $30\,\%\cdot\frac{9}{10}=27\,\%$ of
the labels are actually wrong. With \emph{asymmetric noise}, labels are
flipped to specific similar classes, e.g.\ truck to automobile or cat to dog,
which mimics real annotator confusions.

\paragraph{Memorization.}
A deep network has enough capacity to fit any labeling of its training set,
even random labels. With noisy labels, perfect training accuracy therefore
means that the network has also learned the wrong labels. This is called
\emph{memorization}. Figure~\ref{fig:memorization} shows how it unfolds with
the ordinary cross-entropy loss. For clean examples, the fraction predicted
correctly rises and stays high. For the wrongly labeled examples, the
prediction is compared with their \emph{true} class. Surprisingly, in the
first epochs the network predicts the true class for most of them, although
it was never told that class: it has learned the general pattern of each
class from the clean majority and applies it to the mislabeled examples too.
Later, the fraction predicted as the \emph{wrong training label} rises, and
at the end almost all wrong labels are memorized. With early-learning
regularization (next subsection), the curve for the wrong labels stays as
it was after early learning.

\begin{figure}[ht]
  \centering
  \begin{tikzpicture}[
    ax/.style={-{Latex[length=1.6mm]}},
    lbl/.style={font=\scriptsize}
  ]
    \foreach \dx/\ttl/\reg in {0/Cross-entropy/0, 7.2/Early-learning
        regularization/1} {
      \begin{scope}[xshift=\dx cm]
        \draw[ax] (0,0) -- (5.6,0);
        \node[lbl] at (2.5,-0.6) {epoch};
        \draw[ax] (0,0) -- (0,3.3);
        \node[lbl, rotate=90] at (-0.85,1.5) {fraction of examples};
        \foreach \t in {0,60,120}
          \draw (\t/120*5,0) -- ++(0,-0.08) node[below, lbl] {\t};
        \foreach \v in {0,0.5,1}
          \draw (0,\v*3) -- ++(-0.08,0) node[left, lbl] {\v};
        \node[font=\small] at (2.6,3.55) {\ttl};
        \ifnum\reg=0
          \draw[green!55!black, very thick, dashed]
            plot[domain=0:120, samples=60, smooth]
            ({\x/120*5},{3*(0.35+0.42*(1-exp(-\x/5))-0.70/(1+exp(-(\x-60)/12)))});
          \draw[red!75!black, very thick]
            plot[domain=0:120, samples=60, smooth]
            ({\x/120*5},{3*(0.07-0.03*(1-exp(-\x/5))+0.85/(1+exp(-(\x-75)/14)))});
          \node[lbl, green!45!black] at (1.2,2.55) {true class};
          \node[lbl, red!75!black, align=center] at (4.6,1.25)
            {wrong label\\memorized};
        \else
          \draw[green!55!black, very thick, dashed]
            plot[domain=0:120, samples=60, smooth]
            ({\x/120*5},{3*(0.35+0.53*(1-exp(-\x/8)))});
          \draw[red!75!black, very thick]
            plot[domain=0:120, samples=60, smooth]
            ({\x/120*5},{3*(0.07-0.04*(1-exp(-\x/5)))});
          \node[lbl, green!45!black] at (3.6,2.3) {true class};
          \node[lbl, red!75!black] at (3.6,0.4) {wrong label memorized};
        \fi
      \end{scope}
    }
  \end{tikzpicture}
  \caption{Schematic training curves for the wrongly labeled examples of a
  noisy dataset. With cross-entropy, the network first predicts their true
  class (early learning) and later reproduces their wrong training label
  (memorization). Early-learning regularization keeps the early behavior.
  For the clean examples, the fraction predicted correctly rises in both
  cases.}
  \label{fig:memorization}
\end{figure}

\paragraph{Why memorization comes late.}
The cross-entropy for $n$ examples with one-hot training labels
$\mathbf y^{[i]}$ and softmax outputs $\mathbf p^{[i]}$ is
\begin{equation}
  \mathcal L_{\mathrm{CE}}(\Theta) = -\frac1n\sum_{i=1}^{n}\sum_{c=1}^{C}
    y^{[i]}_c \log p^{[i]}_c .
  \label{eq:weak-ce}
\end{equation}
Its gradient with respect to the logit $z^{[i]}_c$ of example $i$ is the
simple residual
\begin{equation}
  \frac{\partial \mathcal L_{\mathrm{CE}}}{\partial z^{[i]}_c}
    = \frac1n\bigl(p^{[i]}_c - y^{[i]}_c\bigr).
  \label{eq:ce-grad}
\end{equation}
Early in training, all examples have large residuals. The clean examples
form the majority and agree with each other, so their gradients add up and
dominate; the network learns the true class structure. Once the clean
examples are fitted, $\mathbf p^{[i]}\approx\mathbf y^{[i]}$ and their
gradients vanish. The wrongly labeled examples are still predicted as their
true class, which disagrees with their training label, so their residuals
stay close to $1$. From this point on, they dominate the total gradient and
the network gradually bends its decision boundary around them.

\paragraph{The simulated experiment.}
This mechanism was shown analytically for a simple case. The inputs come from
two Gaussians with opposite means,
\begin{equation}
  x\sim\mathcal N(+\mathbf v,\sigma^2 I_{p\times p})\ \text{if } \mathbf y^*=(1,0),
  \qquad
  x\sim\mathcal N(-\mathbf v,\sigma^2 I_{p\times p})\ \text{if } \mathbf y^*=(0,1),
\end{equation}
where $\mathbf v$ is a unit vector and $\mathbf y^*$ the true label. The
optimal separator is the hyperplane through the origin perpendicular to
$\mathbf v$. Each training label equals the true label with probability
$1-\Delta$ and is replaced by the other label with probability $\Delta$. A
linear classifier trained by gradient descent on the cross-entropy then
behaves as follows (informal theorem): for any noise rate $\Delta\in(0,1)$,
under further conditions on dimension and sample size, there is an iteration
$T$ such that
\begin{enumerate}
  \item \emph{early learning succeeds}: for $t<T$, the classifier predicts the
    \emph{true} labels of the wrongly labeled examples better than at
    initialization;
  \item the gradients of the \emph{correctly} labeled examples vanish between
    $t=0$ and $t=T$;
  \item \emph{memorization} occurs as $t\to\infty$.
\end{enumerate}
The network output in the early-learning phase therefore contains useful
information about which labels are wrong. The methods below exploit it.
Early stopping would be the naive remedy, but the right stopping time is
unknown without clean validation data, and the model is still immature at
that point.

\subsection{Early-Learning Regularization}

\emph{Early-learning regularization} (ELR) keeps the network close to what it
predicted during early learning. For each training example it maintains a
\emph{target} vector $\mathbf t^{[i]}$ computed from the network's own past
outputs, typically as a running average over epochs,
\begin{equation}
  \mathbf t^{[i]} \leftarrow \beta\,\mathbf t^{[i]} + (1-\beta)\,\mathbf p^{[i]},
  \qquad \beta\in[0,1) ,
\end{equation}
and adds a term that rewards agreement between the current prediction and
the target:
\begin{equation}
  \mathcal L_{\mathrm{ELR}}(\Theta) = \mathcal L_{\mathrm{CE}}(\Theta)
    + \frac{\lambda}{n}\sum_{i=1}^{n}
    \log\Bigl(1-\bigl\langle \mathbf p^{[i]},\mathbf t^{[i]}\bigr\rangle\Bigr).
  \label{eq:elr}
\end{equation}
The dot product $\langle\mathbf p,\mathbf t\rangle=\sum_c p_c t_c$ measures the
similarity of two probability vectors; it is large when both put their mass
on the same class. Maximizing it makes $\log(1-\langle\mathbf p,\mathbf
t\rangle)$ very negative, which lowers the loss. Because the target is
computed from past outputs rather than from the noisy label, the regularizer
pulls every example towards the class the network believed in early, which
for mislabeled examples is mostly the true class. How to estimate the
target best is still a research question; the running average is a simple
and effective choice.

\paragraph{Effect on the gradient.}
Differentiating the regularizer of one example with respect to its logit
$z_c$ (softmax derivative $\partial p_k/\partial z_c=p_k(\delta_{kc}-p_c)$)
gives
\begin{equation}
  \frac{\partial}{\partial z_c}\log\bigl(1-\langle\mathbf p,\mathbf t\rangle\bigr)
  = -\frac{p_c\,\bigl(t_c-\langle\mathbf p,\mathbf t\rangle\bigr)}
    {1-\langle\mathbf p,\mathbf t\rangle}.
  \label{eq:elr-grad}
\end{equation}
Let $c^*$ be the true class and assume the target puts most weight on it, so
$t_{c^*}>\langle\mathbf p,\mathbf t\rangle$. Then the ELR gradient for $c^*$
is negative, and gradient descent increases $z_{c^*}$. Compare this with the
cross-entropy gradient $p_{c^*}-y_{c^*}$ of Equation~\ref{eq:ce-grad}:
\begin{itemize}
  \item For a \emph{clean} example, $y_{c^*}=1$ and the CE gradient
    $p_{c^*}-1$ is negative early and approaches $0$ once the example is
    fitted. The ELR term adds another negative contribution, so both terms
    agree.
  \item For a \emph{wrongly labeled} example, $y_{c^*}=0$ and the CE gradient
    $p_{c^*}$ is positive; it pushes the true class \emph{down}, which is
    memorization. The ELR gradient has the opposite sign and roughly
    cancels it, so the wrong label cannot take over.
\end{itemize}
In the early phase, the CE gradients of the clean labels dominate and drive
learning. Later, when they vanish, ELR neutralizes the wrong labels instead
of letting them dominate. The clean labels thus stay in charge throughout
training.

\paragraph{Worked example.}
A mislabeled image has true class $c^*$, current output $p_{c^*}=0.6$, and
target $\mathbf t$ with $t_{c^*}=0.9$; assume $\langle\mathbf p,\mathbf
t\rangle=0.56$. The CE gradient for $z_{c^*}$ is $+0.6$ (push down). The ELR
gradient is $-0.6\cdot(0.9-0.56)/(1-0.56)\approx-0.46$ per unit of $\lambda$.
With $\lambda=3$, the regularizer contributes $\approx-1.39$, and the sum is
negative: the true class is strengthened despite the wrong label. Note the
factor $1/(1-\langle\mathbf p,\mathbf t\rangle)$: the closer the prediction
already is to the target, the stronger the regularizer holds it there.

\subsection{Negative Learning for Noisy Labels}

The standard training signal, called \emph{positive learning} (PL) in this
context, tells the network ``this input belongs to this label.'' If the label
is wrong, the statement is wrong, and the network learns a wrong decision
boundary. \emph{Negative learning} (NL) uses a weaker but much safer
statement: ``this input does \emph{not} belong to this other label.'' Take an
image of a dog labeled ``car'' by mistake. PL trains ``this is a car,'' which
is false. NL picks a random class other than the given label, say ``bird,''
and trains ``this is not a bird,'' which is true.

The two losses for $C$ classes are
\begin{align}
  \text{PL:}\quad \mathcal L(f,y) &= -\sum_{k=1}^{C} y_k\log p_k ,
  \label{eq:pl}\\
  \text{NL:}\quad \mathcal L(f,\bar y) &= -\sum_{k=1}^{C} \bar y_k\log(1-p_k) ,
  \label{eq:nl}
\end{align}
where $\mathbf y$ is the one-hot given label and $\bar{\mathbf y}$ the one-hot
\emph{complementary label}. PL pushes $p_y$ towards $1$; NL pushes
$p_{\bar y}$ towards $0$. The complementary label is drawn anew in every
iteration, uniformly from all classes except the given one:
\begin{equation}
  \bar y \sim \mathcal U\bigl(\{1,\dots,C\}\setminus\{y\}\bigr).
\end{equation}

\paragraph{Why NL is robust.}
A complementary label is wrong only if it hits the true class. For a clean
example this is impossible, because the true class is the excluded label.
For a mislabeled example, the true class is one of the $C-1$ candidates, so
the chance is $1/(C-1)$. On CIFAR-10 ($C=10$) with $27\,\%$ wrong labels:
\begin{equation*}
  P(\text{PL statement wrong}) = 0.27,
  \qquad
  P(\text{NL statement wrong}) = 0.27\cdot\tfrac19 = 0.03 .
\end{equation*}
NL thus reduces the rate of misleading training signals by a factor $C-1$.
The price is less information per step: ``not a bird'' excludes one of ten
classes, while ``a car'' fixes the class. NL converges more slowly and on
its own reaches a lower final accuracy than PL on clean data.

\paragraph{Observed behavior.}
With $30\,\%$ symmetric-inclusive noise on CIFAR-10, a CNN trained with PL
reaches almost $100\,\%$ accuracy on the noisy training labels, while its
accuracy on the clean test set peaks early and then drops as it memorizes;
the test loss rises again. The same CNN trained with NL fits the noisy
training labels much less, but its test accuracy keeps increasing and ends
clearly above PL. The difference is visible in the confidence $p_y$ that each
network assigns to the \emph{given} label of each training image
(Figure~\ref{fig:nl-hist}). After PL, both clean and wrongly labeled images
have $p_y$ near $1$: the wrong labels are treated as correct. After NL, the
wrongly labeled images have $p_y$ near $0$, while the clean ones spread over
higher values. The given label is no longer believed where it is wrong, so
$p_y$ becomes a noise detector.

\begin{figure}[ht]
  \centering
  \begin{tikzpicture}[
    lbl/.style={font=\scriptsize}
  ]
    \foreach \dx/\ttl/\clean/\noisy in {
        0/PL/{0,0,0,0,0,0,0.02,0.05,0.12,0.25,0.4,1.0}/{0,0,0,0,0,0,0.03,0.08,0.14,0.17,0.12,0.06},
        4.6/NL/{0.01,0.04,0.06,0.08,0.11,0.14,0.2,0.3,0.45,0.3,0.05,0.0}/{0.9,0.25,0.06,0.02,0.01,0,0,0,0,0,0,0},
        9.2/{SelNL\,$\to$\,SelPL}/{0,0,0,0,0,0,0,0,0,0,0.06,1.0}/{0.8,0.02,0,0,0,0,0,0,0,0,0,0}} {
      \begin{scope}[xshift=\dx cm]
        \foreach [count=\i from 0] \h in \clean
          \fill[blue!45, opacity=0.8] (\i*0.3,0) rectangle ++(0.3,\h*2);
        \foreach [count=\i from 0] \h in \noisy
          \fill[orange!80, opacity=0.8] (\i*0.3,0) rectangle ++(0.3,\h*2);
        \draw[-{Latex[length=1.5mm]}] (0,0) -- (3.8,0);
        \draw[-{Latex[length=1.5mm]}] (0,0) -- (0,2.3);
        \node[lbl, below] at (0,0) {0};
        \node[lbl, below] at (3.6,0) {1};
        \node[lbl, below] at (1.8,-0.05) {$p_y$};
        \node[font=\small] at (1.8,2.55) {\ttl};
      \end{scope}
    }
    \fill[blue!45] (3.2,-1.0) rectangle ++(0.3,0.2);
    \node[lbl, anchor=west] at (3.55,-0.9) {correct label};
    \fill[orange!80] (6.2,-1.0) rectangle ++(0.3,0.2);
    \node[lbl, anchor=west] at (6.55,-0.9) {wrong label};
  \end{tikzpicture}
  \caption{Schematic histograms of the confidence $p_y$ in the given label
  for all training images (CIFAR-10, $30\,\%$ symmetric-inclusive noise).
  After PL, wrong labels look as confident as correct ones. After NL, they
  collect near $0$. Selective NL and PL sharpen the separation until clean
  and noisy data are almost perfectly split.}
  \label{fig:nl-hist}
\end{figure}

\paragraph{Selective NL and selective PL.}
The separation after NL is good but not perfect, and plain NL learns slowly.
Two selective stages refine it, together called SelNLPL:
\begin{itemize}
  \item \emph{Selective negative learning} (SelNL) continues NL, but only on
    examples whose confidence in the given label exceeds $1/C$, i.e.\ is
    higher than uniform guessing. Examples with lower confidence are likely
    mislabeled and are excluded, so they no longer disturb training. After
    SelNL, the clean examples move towards $p_y\approx1$.
  \item \emph{Selective positive learning} (SelPL) then switches to PL, the
    faster and more informative loss, but only for examples with $p_y>\gamma$,
    with $\gamma=0.5$ in practice. These are almost exclusively clean
    examples, so PL can now be used without learning wrong labels. After
    SelPL, nearly all clean data have confidence close to $1$, and the noisy
    data stay near $0$.
\end{itemize}
Table~\ref{tab:selnlpl} lists the three stages. Each runs for a number of
epochs $T$ on the same network $f(x;\theta)$.

\begin{table}[ht]
  \centering
  \begin{tabular}{@{}p{1.9cm}p{4.3cm}p{2.2cm}p{5.4cm}@{}}
    \toprule
    Stage & Examples used in a batch & Loss & Purpose \\
    \midrule
    1.\ NL & all & NL, Eq.~\ref{eq:nl} & robust start; separate clean from
      noisy \\
    2.\ SelNL & only if $p_y>1/C$ & NL, Eq.~\ref{eq:nl} & ignore likely noisy
      examples \\
    3.\ SelPL & only if $p_y>\gamma$ ($\gamma=0.5$) & PL, Eq.~\ref{eq:pl} &
      fast, precise learning on likely clean data \\
    \bottomrule
  \end{tabular}
  \caption{The SelNLPL procedure. The selection threshold rises from stage
  to stage, and the loss switches from the safe negative loss to the
  informative positive loss once the data are filtered.}
  \label{tab:selnlpl}
\end{table}

\paragraph{Pseudo-labeling the noisy part.}
After SelNLPL, the training data are split into a \emph{clean} set (high
$p_y$) and a \emph{noisy} set (low $p_y$). The noisy images are still useful
images, only their labels are not. Treating them as unlabeled data turns the
problem into semi-supervised learning, and \emph{pseudo-labeling} closes the
loop (Figure~\ref{fig:pseudo-label}):
\begin{enumerate}
  \item Divide the training data into clean and noisy data using the CNN
    trained with SelNLPL.
  \item Train a freshly initialized CNN on the clean data only.
  \item Relabel the noisy data with the \emph{soft} output of this CNN, i.e.\
    the full probability vector instead of a one-hot label, which keeps the
    network's uncertainty.
  \item Train another freshly initialized CNN on the clean data plus the
    relabeled data.
\end{enumerate}

\begin{figure}[ht]
  \centering
  \begin{tikzpicture}[
    box/.style={draw, rounded corners=2pt, minimum height=0.9cm,
      minimum width=1.7cm, align=center, font=\scriptsize},
    dat/.style={draw, minimum height=0.6cm, minimum width=1.4cm,
      align=center, font=\scriptsize},
    arr/.style={-{Latex[length=2mm]}, thick},
    lbl/.style={font=\scriptsize}
  ]
    \node[dat, fill=yellow!25] (td) at (0,0) {training\\data};
    \node[box] (n1) at (2.3,0) {CNN\\(SelNLPL)};
    \node[dat, fill=orange!30] (nd) at (4.7,0.85) {noisy};
    \node[dat, fill=blue!15] (cd) at (4.7,-0.85) {clean};
    \node[box] (n2) at (7.0,0) {new CNN};
    \node[dat, fill=orange!55] (rd) at (9.3,0.85) {relabeled\\(soft)};
    \node[dat, fill=blue!15] (cd2) at (9.3,-0.85) {clean};
    \node[box] (n3) at (11.6,0) {new CNN\\(final)};
    \draw[arr] (td) -- (n1);
    \draw[arr] (n1.east) -- (nd.west);
    \draw[arr] (n1.east) -- (cd.west);
    \draw[arr] (cd.east) -- (n2.west);
    \draw[arr] (nd.east) -- (n2.west);
    \draw[arr] (n2.east) -- (rd.west);
    \draw[arr] (rd.east) -- (n3.west);
    \draw[arr] (cd2.east) -- (n3.west);
    \draw[decorate, decoration={brace, mirror, amplitude=4pt}]
      (-0.7,-1.5) -- (5.4,-1.5) node[midway, below=4pt, lbl] {(1) split};
    \draw[decorate, decoration={brace, mirror, amplitude=4pt}]
      (5.6,-1.5) -- (10.0,-1.5) node[midway, below=4pt, lbl]
      {(2)--(3) train on clean, relabel noisy};
    \draw[decorate, decoration={brace, mirror, amplitude=4pt}]
      (10.2,-1.5) -- (12.5,-1.5) node[midway, below=4pt, lbl] {(4) retrain};
  \end{tikzpicture}
  \caption{Pseudo-labeling after SelNLPL. The noisy labels are discarded
  and replaced by soft labels predicted by a network trained on the clean
  part; the final network learns from all images.}
  \label{fig:pseudo-label}
\end{figure}

\paragraph{Results.}
Table~\ref{tab:nl-results} shows test accuracies on CIFAR-10 and CIFAR-100
(ResNet-34). Plain cross-entropy degrades quickly as symmetric noise grows;
the full method stays far above it, e.g.\ $88.3\,\%$ instead of $74.1\,\%$ on
CIFAR-10 with $60\,\%$ noise. With asymmetric noise on CIFAR-100, the
\emph{forward correction} method (Forward $T$) is slightly better. That
method, however, uses the noise transition matrix $T$, i.e.\ knowledge of
which class is mainly confused with which other class. This is information
the negative learning approach does not need.

\begin{table}[ht]
  \centering
  \begin{tabular}{@{}llcccccc@{}}
    \toprule
    & & \multicolumn{3}{c}{symmetric noise} & \multicolumn{3}{c}{asymmetric
      noise} \\
    \cmidrule(lr){3-5}\cmidrule(l){6-8}
    Dataset & Method & 20\,\% & 40\,\% & 60\,\% & 10\,\% & 20\,\% & 40\,\% \\
    \midrule
    CIFAR-10 & CE & 86.98 & 81.88 & 74.14 & 90.69 & 88.59 & 80.11 \\
      & Forward $T$ & 88.63 & 85.07 & 79.12 & 91.32 & 90.35 & 88.12 \\
      & SelNLPL + pseudo-labels & \textbf{94.23} & \textbf{92.43} &
        \textbf{88.32} & \textbf{94.57} & \textbf{93.35} & \textbf{89.86} \\
    \midrule
    CIFAR-100 & CE & 58.72 & 48.20 & 37.41 & 66.54 & 59.20 & 42.74 \\
      & Forward $T$ & 63.16 & 54.65 & 44.62 & \textbf{71.05} & \textbf{71.08}
        & \textbf{70.82} \\
      & SelNLPL + pseudo-labels & \textbf{71.52} & \textbf{66.39} &
        \textbf{56.51} & 70.35 & 63.12 & 45.70 \\
    \bottomrule
  \end{tabular}
  \caption{Test accuracy (\%) under label noise, excerpt from the NLNL paper.
  Negative learning wins clearly under symmetric noise. Under strong
  asymmetric noise on CIFAR-100, forward correction wins, but it assumes a
  known confusion matrix.}
  \label{tab:nl-results}
\end{table}

\subsection{How Much Does Weak Supervision Really Gain?}

Weakly supervised learning is an active research area, and a critical look
at published results shows that the gains are often smaller than claimed.
Many WSL methods silently rely on some \emph{clean} data, for example a
clean validation set to choose hyperparameters or the stopping epoch. If
these clean samples are available anyway, a natural baseline is to train a
standard model, or fine-tune a pretrained one, directly on them. Such
experiments showed that a few clean labels per class often match or beat
WSL methods trained on thousands of weakly labeled examples. Three
recommendations follow for evaluating any WSL approach:
\begin{itemize}
  \item Report how much the method relies on clean data, including clean
    validation data for model selection.
  \item Report how many cleanly annotated samples a few-shot approach needs
    to reach the same performance. If thousands of weakly labeled samples
    are worth only a handful of clean ones, WSL may not be the right
    choice.
  \item If the method needs extra clean data, include standard non-WSL
    methods trained on that clean data in the comparison.
\end{itemize}

\noindent
Table~\ref{tab:weak-summary} collects the methods of this section with
their key idea and assumption. The assumption column is the practical
checklist: a method only helps if its assumption holds for the data at hand.

\begin{table}[ht]
  \centering
  \begin{tabular}{@{}p{3.7cm}p{5.2cm}p{4.9cm}@{}}
    \toprule
    Method & Key idea & Assumption / requirement \\
    \midrule
    Active learning & ask an oracle for the most useful labels & an expert is
      available; query cost matters \\
    Semi-supervised & unlabeled data shape the clusters; labels
      name them & cluster assumption holds \\
    Transductive & label the given unlabeled set directly & test data
      known in advance \\
    Multi-instance & bag positive iff one instance positive; max
      pooling & bag-level labels; negative bags fully negative \\
    ELR & pull predictions towards early-training targets &
      early learning finds the true classes \\
    Negative learning & train ``not this class'' with random complementary
      labels & noise rate moderate; $C$ not too small \\
    SelNLPL & filter by confidence, then PL on clean data and
      relabel the rest & NL confidence separates clean from noisy \\
    \bottomrule
  \end{tabular}
  \caption{Summary of weakly supervised methods. Each exploits a specific
  structure in weak labels and fails when its assumption is violated.}
  \label{tab:weak-summary}
\end{table}
